\section{Virtual degree and the universal class on the same family}\label{sec:comparison}

Let $X\to S$ be the smooth algebraic family of threefolds obtained in
Theorem~\ref{thm:separate}. At a transverse solution the image is a
smooth embedded curve $D\subset X_s$ of multiplicity one. Its normal
bundle $N_D$ has degree $2h-2$, where $h$ is the genus of $D$.
The deformations over $S$ constructed in \S\ref{sec:deformations}
surject onto the obstruction space $H^1(D,N_D)$. Indeed, the
deformations with the projection cover fixed already span the vertical
obstruction space, and the cover itself is unobstructed. The quotient
by deformations of the cover gives the obstruction space of the
embedded map. We first compute the contribution from all connected
maps with image cycle $[D]$.

\subsection{The relative virtual class}

Let $\mathcal H$ be an open neighborhood of $[D]$ in the relative Hilbert
scheme of smooth connected proper curves in $X/S$. It can be taken
inside the Hilbert scheme of a projective completion by restricting
to curves contained in $X$. Shrink $\mathcal H$ so that the normal obstruction
map remains surjective. Write $p:\mathcal C\to \mathcal H$ for the universal
curve and $N$ for its relative normal bundle in $X\times_S \mathcal H$.
Let $\mathcal M_k(X/S)$ be the stack of connected unmarked stable
maps of arithmetic genus $k$ to fibers of $X\to S$. We use its
vertical stable-map obstruction theory over the stack of prestable
curves times $S$, with the fundamental class of the smooth parameter
space. This gives a virtual class on the quasi-projective target
open of \S\ref{sec:deformations}.

\begin{proposition}\label{prop:moving-embedded-curve}
The scheme $\mathcal H$ is smooth of dimension $\dim S$. For $k\ge h$, let
\[
 q:\mathcal D=\overline{\mathcal M}_{k}(\mathcal C/\mathcal H,1)\longrightarrow \mathcal H
\]
be the proper stack of degree-one maps to its fibers. Composition
with the universal embedding identifies $\mathcal D$ with the open
part of $\mathcal M_k(X/S)$ over the multiplicity-one images
parametrized by $\mathcal H$. The virtual class of this open part, pushed
forward by $q$, is $b_{k,h}[\mathcal H]$, where
\begin{equation}\label{eq:embedded-degree-one}
 b_{k,h}=[u^{2k-2}]\bigl(2\sin(u/2)\bigr)^{2h-2}.
\end{equation}
For $k<h$ the inverse image of this image locus is empty.
\end{proposition}

\begin{proof}
For a smooth embedding the tangent complex relative to $S$ is
$Rp_*N$. This follows by canceling the tangent of the source in the
deformation complex of the map. The absolute obstruction at $[D]$
is the cokernel of the map $T_sS\to H^1(D,N_D)$. Its vanishing
gives smoothness, and $\chi(D,N_D)=0$ gives the asserted dimension.
In particular the derived space of smooth embedded maps over the
varying target is classical here. Its absolute tangent is $T_\mathcal H$.

Write $\pi:\mathcal F\to\mathcal D$ for the universal source,
$a:\mathcal F\to\mathcal C$ for the universal map, and
$\iota:\mathcal D\to\mathcal M_k(X/S)$ for composition. The cone of
the induced morphism of virtual tangent complexes is
\begin{equation}\label{eq:comparison-normal-cone}
 \operatorname{Cone}\bigl(q^*Rp_*N\longrightarrow R\pi_*a^*N\bigr).
\end{equation}
To see the arrow, work first over the stack of prestable sources.
With the target curve fixed the vertical tangent is
$R\pi_*a^*T_{\mathcal C/\mathcal H}$, and with the threefold fixed it is
$R\pi_*a^*(T_{X/S}|_{\mathcal C})$. Their quotient is
$R\pi_*a^*N$. Allowing the embedded curve to move adds the
triangle for $T_\mathcal H\to T_S|_\mathcal H$, whose fiber is $Rp_*N$.
A normal displacement of the embedded curve acts on the composite
by pullback under $a$. The two tangent triangles therefore give
\eqref{eq:comparison-normal-cone}. Adding the common motion of the
prestable source leaves the cone unchanged. Composition with the
universal embedding is a morphism of the derived mapping stacks.
Upon truncation it induces the maps of stable-map obstruction
theories, including their maps to the intrinsic cotangent complexes
\cite[\S\S1, 2.2]{STV}. The construction is
relative to the parameter scheme and local in the target along the
proper source.

At a geometric point $a:F\to D$ of $\mathcal D$, the normalization
of the unique positive component maps isomorphically to $D$.
It has no self-node, since two distinct points on the normalization
have distinct images. Each connected contracted subcurve attaches
to this component at just one point, for the same reason. Functions
on a contracted connected fiber are constant and take the value at
the attachment. Hence $a_*\mathcal O_F=\mathcal O_D$ and
$R^1a_*\mathcal O_F$ is a sheaf of finite length $k-h$, by the
Euler characteristics. In particular $k\ge h$. The unit
$\mathcal O_D\to Ra_*\mathcal O_F$ and the projection formula
identify the fiber of \eqref{eq:comparison-normal-cone} with
\[
 H^0(D,N_D\otimes R^1a_*\mathcal O_F)[-1].
\]
Both terms in \eqref{eq:comparison-normal-cone} are perfect and
commute with derived base change along the flat universal curves.
A perfect complex whose geometric fibers are concentrated in one
degree with constant rank is a vector bundle in that degree.
Thus \eqref{eq:comparison-normal-cone} is $W[-1]$ for a bundle
$W$ of rank $2(k-h)$ on $\mathcal D$.

The relative tangent complex of the derived composition morphism
restricts to $W[-2]$ on the classical truncation. For an ordinary
square-zero extension the relative infinitesimal automorphism,
tangent and obstruction groups in degrees $-1,0,1$ consequently
vanish. The classical morphism $\iota$ is formally \'etale and locally
of finite presentation, hence \'etale. Each image determines its
point of $\mathcal H$ uniquely, and the stabilizers on both sides consist of
the same automorphisms of the map from $F$. Thus $\iota$ is a
representable \'etale monomorphism, and identifies $\mathcal D$
with the asserted open substack.

On this common classical stack the vector bundle stack of the
obstruction theory of $\mathcal D$ embeds in the one for
$\mathcal M_k(X/S)$ with quotient $W$. By compatibility the
embedding of the intrinsic normal cone for the latter theory factors
through its embedding for the former. The zero-section Gysin maps
therefore give \cite[\S\S4--5]{BF}
\begin{equation}\label{eq:comparison-excess}
 \iota^*[\mathcal M_k(X/S)]^{\rm vir}
       =c_{2(k-h)}(W)\cap[\mathcal D]^{\rm vir}.
\end{equation}
Here the class on the right uses the smooth family of target curves
over $\mathcal H$. The equality holds on the stacks, including the stabilizers
of all contracted subcurves.

The right side of \eqref{eq:comparison-excess} has dimension
$\dim \mathcal H$. Its proper pushforward is a multiple of $[\mathcal H]$ on each
connected component. To compute the multiple, intersect with a point
$[D]\in \mathcal H$. Refined base change gives the usual virtual class of
$\overline{\mathcal M}_k(D,1)$ \cite[Proposition~7.2]{BF}.
On this fiber the class of $R\Gamma(D,N_D)$ in $K$-theory vanishes
because it is a constant class of rank zero. Consequently the
coefficient is
\[
 \int_{[\overline{\mathcal M}_{k}(D,1)]^{\rm vir}}
             c_{2(k-h)}(-R\pi_*a^*N_D).
\]
In rational cohomology Grothendieck--Riemann--Roch on the universal
source shows that this expression depends only on the rank and
degree of $N_D$. We may use $A\oplus K_DA^{-1}$ for a degree-$(h-1)$
line $A$ outside the theta divisor. Both summands are acyclic,
and a section of their pullback to $F$ is determined by its
restriction to the positive component. The integral is thus the
local degree-one contribution computed by
Pandharipande \cite[Theorem~2]{Pand}. Multiplying his series in
the excess genus by $u^{2h-2}$ gives \eqref{eq:embedded-degree-one}.
\end{proof}

\subsection{The local class in families of cycles}

Let $\sigma:B\to S$ be the transverse smooth section of
Theorem~\ref{thm:separate}. Give its normal bundle the orientation
of the vertical tangent of $S\to B$, and give the tests their complex
orientations. At a solution with image $D$, let $\epsilon(D)$ be
the local intersection sign on $\mathcal H$ of the parameter graph and the
tests. The total real codimension is $2\dim_{\mathbb C} \mathcal H$.
The parameter graph is transverse before the tests are imposed, so
its inverse image in $\mathcal H$ is the regular space of embedded curves in
the family over $B$. Its transverse-fiber-product orientation agrees
with the determinant orientation of the parametrized map equation,
by the tangent exact sequence for the graph. Interchanging the
normal to the graph and the tests gives no sign, as their real ranks
are even. Proposition~\ref{prop:moving-embedded-curve} and the
projection formula now give $\epsilon(D)b_{k,h}$ for the
unnormalized algebraic contribution at this solution.

\begin{lemma}\label{lem:signed-embedded-cycle}
In the smooth family colimit the local unnormalized test class at
this solution equals $\epsilon(D)\Phi(x_{h,1})$.
\end{lemma}

\begin{proof}
Embedded connected curves of multiplicity one form an open subset
of the space of cycles, with their embedded-map topology
\cite[Proposition~2.21]{DIW}. Thus the transverse solution can be
excised in cycle space. To identify its class we first deform the
complex structure of the marked target germ so that the curve
becomes regular in that target, keeping the curve and its test
values fixed.

Write $V=H^0(D,N_D)$ and $U=H^0(D,K_D\otimes N_D^\vee)$, both
of dimension $l$ by Riemann--Roch. If $l>0$, choose points
$p_1,\dots,p_l$ on unramified sheets over distinct points of the
base curve away from the test fibers. Choose covectors $\lambda_i$
and vectors $v_i$ in the respective normal fibers so that both maps
\[
 V\longrightarrow\mathbb C^l,\quad
       s\longmapsto(\lambda_i s(p_i))_i,\qquad
 U\longrightarrow\mathbb C^l,\quad
       \eta\longmapsto(\eta(p_i)(v_i)/dx)_i
\]
are isomorphisms. Here $x$ is a local coordinate on the base.
Inductively, nonzero sections in the two remaining evaluation
kernels are simultaneously nonzero on some unramified sheet. At
such a point choose the vector and covector to reduce each kernel
dimension by one. Extend these data in local frames. Multiply the
endomorphisms $v_i\lambda_i$ by smooth $(0,1)$-forms supported in
small disks about $p_i$, normalized to pair to one with $dx$, and
denote their sum by $A$. Under Serre duality the matrix of
$J:V\to H^1(D,N_D)$, $s\mapsto[As]$, tends to the product of the
two evaluation matrices, one transposed, as the disks shrink. Hence
it is invertible for a fixed choice of the disks. Splitting off the
invertible complement of kernel and cokernel in $\bar\partial_{N_D}$
gives the matrix $tJ+O(t^2)$ on these finite-dimensional
spaces for $\bar\partial_{N_D}+tA$. This operator is invertible for all
sufficiently small $t>0$.

Over the chosen base disks take disjoint holomorphic submersion
charts about the sheets. Use the base coordinate $x$ and vertical
coordinates $y$ in which the marked sheet is $y=y_i(x)$.
Over the support of the perturbation a sufficiently small target
tube meets only these charts. Write $A=A_i(x)d\bar x$ on the
chosen sheet and use zero on the other sheets. Replace its local
graph equation by
\[
        \partial_{\bar x}y=-tA_i(x)(y-y_i(x)).
\]
The coefficient vanishes near the boundary of the base disk, so
the change glues on the tube. Its holomorphic dependence on $y$
makes the horizontal $(0,1)$ vector field commute with the two
vertical $(0,1)$ coordinate fields. Hence the resulting almost
complex structure is integrable. Along the marked sheet the added vector
field vanishes and its normal linearization is the prescribed
change of Dolbeault operator. Extend the change smoothly over a
parameter neighborhood in $B$, with a parameter cutoff. All framed
test fibers are unchanged. For $l=0$ the fixed target was already
regular.

At $t=0$ the differential of the parametrized map equation together
with the tests is invertible in an unparametrized slice. The implicit
function theorem therefore isolates the retained solution on a small
closed interval of $t$. By the open embedding into the relative
cycle space its local Thom class extends by zero over the interval,
with the same sign. At the regular endpoint the local family of
curves projects diffeomorphically to $B$, preserving orientation
since the fixed-fiber normal operator is an invertible complex
operator. The test class is now the base Gysin image of the unit of
the isolated curve, multiplied by $\epsilon(D)$.

Deform the holomorphic embedding to its normal bundle. The curve
stays regular along this deformation, with the same normal
Dolbeault operator up to isomorphism at the positive parameters.
The implicit function theorem and embedded-cycle openness isolate
its compact section along the full interval. At the normal endpoint
a compact degree-one cycle is a section, since the vector-space
fibers have no compact positive curves. Acyclicity makes the zero
section the only such cycle. Its unit is the empty-tuple
representative of the multiplicity-one class of genus $h$ and Chern
number zero \cite[Definition~4.7]{Pardon}. The two local
deformations and the base Gysin map identify the asserted class
\cite[\S\S2--3]{WueProduct}.
\end{proof}

\subsection{Extraction by source genus}

\begin{proof}[Proof of Theorem~\ref{thm:main}]
Fix $K$ and use Theorem~\ref{thm:separate} for all connected source
genera through $K$. For each $h\le K$ let $a_h^{(K)}$ be the sum
of $\epsilon(D)$ over the finitely many connected embedded solutions
of genus $h$, divided by $W_{g,m}$. The parameter graph and the
tests form a single compactly supported intersection class on the
algebraic family, isotopic to the graph and tests at the original
algebraic section. Hence the pairing computed in
\S\ref{sec:weights} is the sum of the local pairings in
Proposition~\ref{prop:moving-embedded-curve}. For every $k\le K$
this gives
\begin{equation}\label{eq:connected-algebraic-extraction}
             N_{k,m,j}=\sum_{h\le k} a_h^{(K)} b_{k,h}.
\end{equation}
Here and below a smaller source genus cannot map to a
multiplicity-one image of larger genus.

In the smooth family over the moving graph the same test classes
give a compactly supported homotopy on cycles. For fixed $k\le K$
the image-cycle map from connected genus-$k$ stable maps is proper
\cite[\S5]{WueProduct}. Outside neighborhoods of the finitely many
embedded solutions through genus $K$, the compact constrained set
misses its closed image. The pullback of that part of the test
class is therefore zero. Over each retained solution
Lemma~\ref{lem:signed-embedded-cycle} applies before this
pullback, including when the inverse image contains contracted
components. Finally, the product terms subtracted in the logarithm
require positive images in disjoint targets and have zero connected
pullback. Thus
\begin{equation}\label{eq:connected-universal-extraction}
          T_k(P_{m,j})=\sum_{h\le k} a_h^{(K)}T_k(x_{h,1}).
\end{equation}

In the range of weights of the theorem, the primitive-division
lemma gives a finite expansion $P_{m,j}=\sum_h a_{h,j}x_{h,1}$
in the complex cycle algebra, as in \S\ref{sec:weights}. Since
$T_k(x_{h,1})=0$ for $h>k$ and $T_k(x_{k,1})\ne0$
\cite[\S\S5--6]{WueProduct}, induction on $k$ in
\eqref{eq:connected-universal-extraction} identifies
$a_{h,j}=a_h^{(K)}$ for $h\le K$. Equation
\eqref{eq:connected-algebraic-extraction} therefore holds with the
coefficients of the complex universal class for every $k$, since
$K$ was arbitrary. Substitution of \eqref{eq:embedded-degree-one}
gives the asserted sine expansion. The fixed-genus residue of
\S\ref{sec:weights} vanishes for $k<1+m(g-1)+|j|$, and
$b_{h,h}=1$. The same induction in
\eqref{eq:connected-algebraic-extraction} gives the stated lower
bound on the genus of a nonzero coefficient.
\end{proof}
